\documentclass[submission,copyright,creativecommons]{eptcs}
\providecommand{\event}{SOS 2007} % Name of the event you are submitting to

\usepackage{iftex}

\ifpdf
  \usepackage{underscore}         % Only needed if you use pdflatex.
  \usepackage[T1]{fontenc}        % Recommended with pdflatex
\else
  \usepackage{breakurl}           % Not needed if you use pdflatex only.
\fi

\title{An EPTCS Example with BibLaTeX Bibliography}
\author{Rob van Glabbeek
\institute{NICTA\\ Sydney, Australia}
\institute{School of Computer Science and Engineering\\
University of New South Wales\thanks{A fine university.}\\
Sydney, Australia}
\email{rvg@cs.stanford.edu}
\and
Co Author \qquad\qquad Yet S. Else
\institute{Stanford University\\
California, USA}
\email{\quad is@gmail.com \quad\qquad somebody@else.org}
}

\newcommand{\titlerunning}{A Longtitled Paper}
\newcommand{\authorrunning}{R.J. van Glabbeek, C. Author \& Y.S. Else}

\hypersetup{
  bookmarksnumbered,
  pdftitle    = {\titlerunning},
  pdfauthor   = {\authorrunning},
  pdfsubject  = {EPTCS},               % Consider adding a more appropriate subject or description
  % pdfkeywords = {keyword1, keyword2} % Uncomment and enter keywords specific to your paper
}

% For the biblatex bibliography:
% * specify `bibstyle=eptcs-numeric`
% * for automatically abbreviating given names by initials, use `giveninits` or `giveninits=true`
\usepackage[bibstyle=eptcs-numeric,giveninits]{biblatex}
% and load your bibliography database file with `\addbibresource`
\addbibresource{generic-biblatex.bib}

\begin{document}
\maketitle

\begin{abstract}
    This EPTCS template shows how to use Bib{\LaTeX} instead of Bib{\TeX}.
    Regarding other matters, see the main example file.
\end{abstract}

Currently, with Bib{\LaTeX}, only a numeric bibliography style is supported.

\section{The Preamble}

In the preamble, load the \texttt{biblatex} package with the option \texttt{bibstyle=eptcs-numeric}.
Optionally, the option \texttt{giveninits} can be given, in which case the given names of authors and editors are automatically abbreviated to initials.
For example, the command for loading biblatex could look as follows:
\begin{center}\ttfamily
    {\textbackslash}usepackage[bibstyle=eptcs-numeric,giveninits]\{biblatex\}
\end{center}
Next, specify the bibliography database file with \texttt{{\textbackslash}addbibresource}.

\section{The Bibliography}

In the document, you can use the cite commands provided by biblatex, e.g.\@ \texttt{{\textbackslash}cite} and \texttt{{\textbackslash}textcite}.
At the point where you want to insert the bibliography, use \texttt{{\textbackslash}printbibliography}.

\section{The Bibliography Database File}

Please follow the biblatex datamodel.
However, note that not all fields documented in the biblatex manual are supported yet, so check the output of a {\LaTeX} compile.

\paragraph{DOIs}

Almost all publishers use digital object identifiers (DOIs) as a
persistent way to locate electronic publications. Prefixing the DOI of
any paper with {\ttfamily https://doi.org/} yields a URI that resolves to the
current location (URL) of the response page\footnote{Nowadays, papers
  that are published electronically tend
  to have a \emph{response page} that lists the title, authors and
  abstract of the paper, and links to the actual manifestations of
  the paper (e.g., as {\ttfamily dvi} or {\ttfamily pdf} file). Sometimes
  publishers charge money to access the paper itself, but the response
  page is always freely available.}
of that paper. When the location of the response page changes (for
instance, through a merge of publishers), the DOI of the paper remains
the same and (through an update by the publisher) the corresponding
URI will then resolve to the new location. For that reason, a reference
ought to contain the DOI of a paper, with a live link to the corresponding
URI, rather than a direct reference or link to the current URL of the
publisher's response page. This is the r\^ole of the biblatex field {\ttfamily doi}.
{\bfseries EPTCS requires the inclusion of a DOI in each cited paper, when available.}

DOIs of papers can often be found through
\url{https://www.crossref.org/guestquery};\footnote{For papers that will appear
  in EPTCS and use \href{https://eptcs.web.cse.unsw.edu.au/eptcs.bst}
  {\ttfamily $\backslash$bibliographystyle$\{$eptcs$\}$} there is no need to
  find DOIs on this website, as EPTCS will look them up for you
  automatically upon submission of the first version of your paper;
  these DOIs can then be incorporated into the final version, together
  with the remaining DOIs that need to be found at DBLP or the publisher's web pages.}
the second method {\itshape Search on article title}, only using the {\bfseries
surname} of the first-listed author, works best.
Other places to find DOIs are DBLP and the response pages for cited
papers (maintained by their publishers).

\paragraph{The fields {\ttfamily eprint} and {\ttfamily url}}

Often an official publication is only available against payment. However,
as a courtesy to readers that do not wish to pay, the authors also
make the paper available free of charge at a repository such as
\url{arXiv.org}. In such a case, it is recommended to also refer and
link to the URL of the response page of the paper in such a
repository.  This can be done using the biblatex fields {\ttfamily eprint} and \texttt{eprinttype}.
The biblatex bibliography style supports, among others, the eprinttypes arxiv and jstor.
If no supported eprint is available, the field {\ttfamily url} can be used.
This field should \textbf{not} be used to duplicate information that is also provided through
{\ttfamily doi} or {\ttfamily eprint}. You can find archival-quality URLs for most recently
published papers in DBLP, but please suppress repetition of DOI or {\ttfamily eprint} information
through {\ttfamily url}. In fact, it is often useful to check your references against DBLP records anyway,
or just find them there in the first place.

\paragraph{References to papers in the same EPTCS volume}

To refer to another paper in the same volume as your own contribution,
use a bibtex entry with
\begin{center}
  {\ttfamily series    = $\{\backslash$thisvolume$\{5\}\}$},
\end{center}
where 5 is the submission number of the paper you want to cite.
You may need to contact the author, volume editors, or EPTCS staff to
find that submission number; it becomes known (and unchangeable)
as soon as the cited paper is first uploaded at EPTCS\@.
Furthermore, omit the fields {\ttfamily publisher} and {\ttfamily volume}.
Then in your main paper, you put something like:

\noindent
{\small \ttfamily $\backslash$providecommand$\{\backslash$thisvolume$\}$[1]$\{$this
  volume of EPTCS, Open Publishing Association$\}$}

\noindent
This acts as a placeholder macro-expansion until EPTCS software adds
something like

\noindent
{\small \ttfamily $\backslash$newcommand$\{\backslash$thisvolume$\}$[1]%
  $\{\{\backslash$eptcs$\}$ 157$\backslash$opa, pp 45--56, doi:\dots$\}$},

\noindent
where the relevant numbers are pulled out of the database at publication time.
Here the newcommand wins from the providecommand, and {\ttfamily \small $\backslash$eptcs}
resp.\ {\ttfamily \small $\backslash$opa} expand to

\noindent
{\small \ttfamily $\backslash$sl Electronic Proceedings in Theoretical Computer Science} \hfill and\\
{\small \ttfamily , Open Publishing Association} \hfill .

\noindent
Hence putting {\small \ttfamily $\backslash$def$\backslash$opa$\{\}$} in
your paper suppresses the addition of a publisher upon expansion of the citation by EPTCS\@.
An optional argument like
\begin{center}
  {\ttfamily series    = $\{\backslash$thisvolume[EPTCS]$\{5\}\}$},
\end{center}
overwrites the value of {\ttfamily \small $\backslash$eptcs}.

\nocite{*}
% Use `\printbibliography` at the point where you want to include the bibliography in your document
\printbibliography
\end{document}
